\chapter{Water accounting and time stepping}\label{ch:c5}
This chapter explains how the two models keep their water accounts consistent, what exactly happens in each time step,
and how runs are started, restarted and repeated.

\begin{plain}
Every day the two models take turns: Raven works out how much water goes down; MODFLOW moves it through the aquifer and
reports what comes back up; Raven then books the returned water. Both keep a running account, and the two accounts are
compared every day. This chapter also explains how to stop a run and carry on later (a restart), and how to run many
versions of a model at once (an ensemble).
\end{plain}

\section{Keeping the books: water balance}\label{ch:balance}
\subsection{MODFLOW's balance}
For every step Raven computes all boundary flows from MODFLOW's own matrix terms ($q=\mathrm{HCOF}\cdot h-\mathrm{RHS}$)
and the storage change from MODFLOW's storage package, and reports the imbalance
\begin{multline}
\varepsilon=Q^{\mathrm{rch}}+Q^{\mathrm{sto}}-Q^{\mathrm{seep}}+Q^{\mathrm{riv}}_{\mathrm{in}}-Q^{\mathrm{riv}}_{\mathrm{out}}-Q^{\mathrm{evt}}
+Q^{\mathrm{wel}}+Q^{\mathrm{ghb}}+Q^{\mathrm{chd}}\\+Q^{\mathrm{swx}}_{\mathrm{in}}-Q^{\mathrm{swx}}_{\mathrm{out}}+Q^{\mathrm{res}}+Q^{\mathrm{sz}}_{\mathrm{in}}-Q^{\mathrm{sz}}_{\mathrm{out}},
\end{multline}
where $Q^{\mathrm{sto}}$ is water released from storage and $Q^{\mathrm{sz}}$ the soil-zone exchange (packages \cmd{SZX\_RAVEN},
\cmd{SZQ\_RAVEN} and \cmd{SZD\_RAVEN}, only with \cmd{:SoilZone PHYSICAL}; the upward flows out of a lower cell through
\cmd{SZW\_RAVEN} are booked with the terms they belong to: the columns' upflow in $Q^{\mathrm{rch}}$, the rise and the
debt repayment below). The residual is one sum for every case that knows the storage flow, the soil-zone terms included
(a steady existing model without \cmd{STO6} has storage 0). It is written in m$^3$/d and, as \emph{MF6 balance error [\%]},
as $100$ times the residual over the \emph{gross flow} through the aquifer: every entry's inflow and outflow apart
(recharge per cell, storage per cell, each well, GHB, specified-head, river, drain, reservoir and evaporation entry, each
soil-zone link net, every entry of an existing model's own packages and its mover), that is MODFLOW's TOTAL IN plus TOTAL
OUT, at least 1~m$^3$/d. This percentage is half MODFLOW's PERCENT DISCREPANCY, which divides by the mean of IN and OUT; the
audit tool's cumulative ratio uses signed package totals, a stricter figure. In the tests of Chapter~\ref{ch:verification}
the worst single-day error is about $5\times10^{-6}$\% and the cumulative error about $10^{-10}$.

\subsection{Raven's balance}
Raven tracks cumulative input $I$ and output $O$ and reports $\mathrm{MB}=\left(S-S_0\right)+\left(O-I\right)$ in mm over
the watershed. The coupling adds:
\begin{itemize}
\item \textbf{to output}: recharge (the end-of-day \cmd{GROUNDWATER} value of coupled HRUs, as Raven already does), and
reservoir seepage (already counted by Raven as a reservoir loss);
\item \textbf{to input}: every volume the coupling adds to Raven in the step --- seepage returned to HRU stores, the net
river exchange applied to reaches, the net surface-store exchange, and with \cmd{:SoilZone PHYSICAL} the soil-zone
exchange applied to the soil layers (rejected recharge plus the net exchange). A net loss (for example a losing river)
is a negative input.
\end{itemize}
With these, Raven's own balance closes exactly as without groundwater (0.43--0.52 mm maximum over 20 years in the Liard
example, the same as uncoupled).

\subsection{The whole system}
The water that leaves Raven as recharge enters MODFLOW (directly or via the delay store); the water MODFLOW sends back
enters Raven. Reservoir seepage enters MODFLOW one day after leaving Raven. Therefore
\begin{equation}
\sum V^{\mathrm{Raven}}_{\mathrm{rch}}=\sum V^{\mathrm{MF6}}_{\mathrm{rch}}+\Delta V^{\mathrm{dl}}+\sum V^{\mathrm{CHD}}_{\mathrm{rch}}+\sum V^{\mathrm{rej}},
\end{equation}
where $V^{\mathrm{rej}}$ is the recharge rejected on saturated ground (zero unless \cmd{:SoilZone PHYSICAL}). With the soil
zone coupled, the volume applied to the soil layers in each step also satisfies
\begin{equation}
V^{\mathrm{sz}}_{\mathrm{applied}}=V^{\mathrm{rej}}+\left(Q^{\mathrm{sz}}_{\mathrm{out}}-Q^{\mathrm{sz}}_{\mathrm{in}}\right)\dt+V^{\mathrm{short}},
\end{equation}
with $V^{\mathrm{short}}$ the shortfall (water MODFLOW took that the soil zone did not hold). The shortfall is 0 whenever the
interface converges; it can appear only on steps whose interface did not converge (most at a low outer-iteration limit).

\paragraph{The soil-zone identities.}\label{sec:accsz} With the unsaturated columns the recharge passes through them, and
an existing model without \cmd{NEWTON} may leave recharge unapplied: the first identity becomes
\begin{equation}
\sum V^{\mathrm{Raven}}_{\mathrm{rch}}=\sum V^{\mathrm{MF6}}_{\mathrm{rch}}+\Delta V^{\mathrm{dl}}+\sum V^{\mathrm{CHD}}_{\mathrm{rch}}+\sum V^{\mathrm{rej}}
+\sum\big(V^{\mathrm{uz}}_{\mathrm{in}}-V^{\mathrm{uz}}_{\mathrm{to\,aq}}\big)+\sum V_{\mathrm{not\ applied}},
\end{equation}
with $V_{\mathrm{not\ applied}}$ the recharge MODFLOW did not apply (signed; Section~\ref{sec:szinterface}, Existing
models). $V^{\mathrm{rej}}$ is negative where a negative recharge was more than the top cell could give: the cut is taken back
out of the soil zone (Section~\ref{sec:negrecharge}).
and the columns themselves close as
$\Delta V^{\mathrm{uz}}=V^{\mathrm{uz}}_{\mathrm{in}}-V^{\mathrm{uz}}_{\mathrm{to\,aq}}-V^{\mathrm{uz}}_{\mathrm{cap}}$ (storage relative to the
hydrostatic profile, Section~\ref{sec:szvadose}), checked from step 1 with the storage at the start that
\cmd{GWModelSummary.txt} gives (``unsaturated storage''). The water applied to the soil zone gains the capillary rise (from the columns
and from the water table), the exchange with specified heads and the lagged-exchange debt:
\begin{equation}
V^{\mathrm{sz}}_{\mathrm{applied}}=V^{\mathrm{rej}}+\left(Q^{\mathrm{sz}}_{\mathrm{out}}-Q^{\mathrm{sz}}_{\mathrm{in}}\right)\dt+V^{\mathrm{short}}
+V^{\mathrm{uz}}_{\mathrm{cap}}+V^{\mathrm{wt}}_{\mathrm{cap}}+\left(Q^{\mathrm{chd}}_{\mathrm{out}}-Q^{\mathrm{chd}}_{\mathrm{in}}\right)\dt+V^{\mathrm{debt}},
\end{equation}
and MODFLOW's balance includes the rise from the water table and the debt repayment as outflows through
\cmd{SZQ\_RAVEN} (or \cmd{SZW\_RAVEN} when they come out of a cell below a dried top cell). The soil zone's own balance,
per HRU over its columns, adds the cut of a negative recharge that the soil zone did not hold, $V^{nh}$:
\begin{equation}
\Delta W\,A/1000=W_{\mathrm{proc}}+V^{\mathrm{rej}}-\textstyle\sum q\,\dt-V^{\mathrm{ovf}}+V^{nh},
\end{equation}
with $\Delta W$ the change of the soil water [mm] over the area $A$, $W_{\mathrm{proc}}$ what Raven's processes moved and
$V^{\mathrm{ovf}}$ the overflow. Each of these terms is a column of \cmd{GWBudget.csv} (Section~\ref{app:outputs}).

The audit tool \cmd{tools\_audit.py} checks these identities, the river, seepage and store bookkeeping, and MODFLOW's
balance for every run; its guard count counts the warnings ``\ldots\ of discharge to Raven withheld'' (without a soil zone
``(MODFLOW unconverged): \ldots\ withheld for MODFLOW's negative residual'', and with \cmd{:SoilZone PHYSICAL}).

\section{A real day of water accounts}
\begin{example}[: 15 June 1996 in the Liard run]
From \cmd{GWBudget.csv} (all in m$^3$/d):
\begin{center}\begin{tabular}{lr}\toprule\hdr
Inflows to the aquifer & \\\midrule
recharge from Raven & 21\,365\,037\\
river leakage & 2\,978\,677\\\midrule\hdr
Outflows from the aquifer & \\\midrule
seepage to Raven & 19\,087\,047\\
discharge to rivers & 1\,331\,456\\\midrule
storage release (negative = stored) & $-$3\,925\,212\\\bottomrule
\end{tabular}\end{center}
Inflows $-$ outflows $+$ release $=24\,343\,714-20\,418\,502-3\,925\,212=0.00003$ m$^3$/d, i.e.\ a relative error of
about $10^{-10}$. Raven recorded the same numbers on its side: 21.4 million m$^3$ of recharge leaving its stores, 19.1 million m$^3$
of seepage entering them, and a net river exchange of $-1.65$ million m$^3$ applied to its reaches (the rivers lost more than
they gained that day). The aquifer stored the difference: a water-table rise of about 1 mm on average.
\end{example}

\section{How to read the balance outputs}
\begin{enumerate}
\item \emph{MF6 balance error [\%]} should stay well below 0.1\% every step; larger values point to a non-converged step or a
too-loose head tolerance (\cmd{:SolverHead\allowbreak Tolerance}, default 0.00001 m).
\item \emph{Raven recharge} and \emph{MF6 recharge} are equal unless a recharge delay or fixed-head cells are used, or
recharge is rejected on saturated ground (\cmd{:SoilZone PHYSICAL}).
\item With \cmd{:SoilZone PHYSICAL}: \emph{soil-zone shortfall} is zero and \emph{interface converged} is 1 on (nearly) every
step; a step that did not converge is also reported in \cmd{Raven\_errors.txt} with the HRU and cell. A shortfall
(water created) can appear only on steps where MODFLOW stopped at a low \cmd{:SolverMaxOuterIterations} (below about
20) without converging, together with a MODFLOW balance error; raise the limit.
\item \emph{returned to Raven} equals \emph{seepage to Raven}.
\item \emph{river loss carried to next step} is normally zero; when not, it is taken from the river the next day.
\item Raven's \cmd{Watershed\allowbreak Storage.csv} MB error should be as small as without groundwater.
\end{enumerate}
The tool \cmd{tools\_audit.py} checks all of these automatically.

\section{The coupled time step, start-up and restarts}\label{ch:timestep}
\subsection{The algorithm}
\begin{figure}[!htbp]\centering
\begin{tikzpicture}[x=1cm,y=1cm,b/.style={draw,rounded corners,font=\scriptsize,align=center,minimum height=8mm,text width=23mm}]
\node[b,fill=teal!12] (a) at (0,0) {Raven HRU\\processes};
\node[b,fill=gray!15] (b) at (2.9,0) {recharge and\\stresses to MF6};
\node[b,fill=violet!12] (c) at (5.8,0) {MODFLOW\\Newton solve};
\node[b,fill=gray!15] (d) at (8.7,0) {seepage, stores,\\river terms back};
\node[b,fill=teal!12] (e) at (11.6,0) {Raven routing\\and balance};
\foreach \u/\v in {a/b,b/c,c/d,d/e} \draw[-{Stealth[length=2mm]},thick] (\u) -- (\v);
\draw[->,thick,gray] (-1.2,-0.8) -- (12.8,-0.8) node[right,font=\scriptsize]{time};
\end{tikzpicture}
\caption{Order of work within one time step.}\label{fig:timeline}
\end{figure}

Algorithm~\ref{alg:step} lists exactly what happens in one Raven time step of a coupled model.

\begin{algorithm}[!htbp]\caption{One coupled time step}\label{alg:step}
\begin{algorithmic}[1]
\State Raven: vertical and lateral processes of all HRUs; coupled HRUs' \cmd{GROUNDWATER} then holds the day's recharge
\State Complete yesterday's line of \cmd{GWBudget.csv} (river and reservoir terms are final only after routing)
\State Recharge: delay reservoirs (Ch.~\ref{ch:recharge}); spread over cells; skip fixed-head cells
\State MODFLOW \cmd{prepare\_time\_step}; set RCH rates
\State Set river stages and conductances, well rates, GHB/CHD heads (Ch.~\ref{ch:rivers},\ref{ch:wells})
\State Set water-table ET rates from PET left unmet by Raven (Ch.~\ref{ch:et})
\State Deliver yesterday's reservoir seepage (Ch.~\ref{ch:reservoirs}); set surface-store levels and limited conductances (Ch.~\ref{ch:stores})
\State \cmd{:SoilZone PHYSICAL}: share Raven's soil changes among the soil-zone columns; connected fractions, rejected recharge, local solves, \cmd{SZX\_RAVEN} boundaries (Ch.~\ref{ch:soilzone})
\State MODFLOW \cmd{prepare\_solve}; \textbf{repeat} \cmd{solve} (with \cmd{:SoilZone PHYSICAL}: then update the soil-zone interface) \textbf{until} converged or the iteration limit; \cmd{finalize\_solve}; \cmd{finalize\_time\_step}
\State Compute every package flow as $\mathrm{HCOF}\cdot h-\mathrm{RHS}$; reservoir local heads for today's routing
\State Surface-store exchange applied to Raven stores; seepage and river flows computed
\If{\cmd{:SoilZone PHYSICAL}} withhold from all discharge to Raven (seepage, gains, soil-zone inflow), in proportion, the soil-zone shortfall the booking would leave and, if the step did not converge, MODFLOW's negative residual; warn
\ElsIf{the step did not converge} withhold MODFLOW's negative residual from all discharge to Raven (seepage, river and store gains), in proportion; warn (first ten) \EndIf
\State \cmd{:SoilZone PHYSICAL}: MODFLOW's soil-zone flows and the rejected recharge applied to the columns; overflow to \cmd{SURFACE\_WATER}
\State Seepage added to Raven stores; balance, heads, connectivity and NetCDF output written
\State Raven: before routing, river gains added to / losses taken from each subbasin's in-catchment inflow
\State Raven routing; losses not covered are taken from the flow entering each reach; reservoirs compute seepage
\State Raven: water balance (recharge out, all returned water in); diagnostics including \cmd{GW\_HEAD}
\end{algorithmic}
\end{algorithm}

\subsection{Starting conditions}\label{sec:init}
Initial heads are chosen in this order: (1) a hotstart block \cmd{:GWHeads} in the \cmd{.rvc} file; (2)
\cmd{:InitialGWHeads} in the \cmd{.rvc} file (\cmd{DEPTH\_BELOW\_SURFACE d} or \cmd{ELEVATION z}); (3) each profile's
\cmd{INITIAL\_HEAD\_DEPTH}. Guessed initial heads can be far from reality; \cmd{:GWInitialization STEADY\_STATE r}
first solves a steady state before Raven's first day, with recharge $r$ [mm/d] over the area covered by coupled HRUs
(cell $i$ receives $r\,A^{\mathrm{cov}}_i/1000$ m$^3$/d, Section~\ref{sec:linkage}) and the initial river stages, well rates and boundary heads. This is skipped on a hotstart. A steady state may not
exist --- for example when wells pump more than the aquifer can supply --- and then the solve does not converge; Raven
warns and starts from the initial heads instead, as if \cmd{:GWInitialization} had not been given. The groundwater model is built after \emph{all} input files,
including the \cmd{.rvc}, have been read.

\subsection{Restarting a run (hotstart)}
At the end of a run Raven writes \cmd{solution.rvc}. Besides Raven's own states it contains \cmd{:GWHeads} (the head in
every cell), \cmd{:GWRechargeStore} (recharge delay stores, if used), \cmd{:GWRiverLossCarry} (river loss carried, if any) and \cmd{:GWReservoirSeepage} (reservoir seepage in
transit, if used), and with \cmd{:SoilZone PHYSICAL} \cmd{:GWSoilZoneColumns} (the water of every soil-zone column), and with unsaturated columns \cmd{:GWVadoseColumns} (\cmd{:GWVadoseClassColumns} with depth classes) and \cmd{:GWVadoseColumnTops} (each column's top and water). Using this file as the \cmd{.rvc} of a new run --- with \cmd{:StartDate} set to the end of the first
run --- continues exactly: in the tests the restarted run matched a continuous run to within $6\times10^{-6}$ in all flows.
Two details make this hold:
\begin{itemize}
\item Without \cmd{:SoilZone PHYSICAL}, seepage is added to its target store (\cmd{:SeepageReturnTo}, by default the
lowest soil layer) with no capacity limit, so that layer can hold more than its capacity for a while. On reading the
\cmd{.rvc}, Raven cuts every state above its maximum back to the maximum; for the seepage target of a coupled HRU it keeps
the value as written, so the restart carries the same water as the continuous run. The same holds for the store of a
\cmd{:SurfaceWaterExchange} (for example \cmd{DEPRESSION} above \cmd{DEP\_MAX}), which groundwater discharge fills with no
capacity limit.
\item The unsaturated columns are restored with the contents as written (\cmd{:GWVadoseColumns},
\cmd{:GWVadoseClassColumns}), floored only at what the method can hold: the air-dry content
$\theta(\psi=-10^5$\,m$)$ for \cmd{RICHARDS} and \cmd{LAG}, zero for \cmd{KINEMATIC} (whose reference content
$\theta_s-S_y$ can lie below air-dry when $S_y$ is close to the porosity). A line is used only for a column with the same
top (within 1~mm, \cmd{:GWVadoseColumnTops}), the same number of nodes and the same depth classes; another column starts at
rest (its reference profile), with a warning. A run without unsaturated columns (\cmd{:SoilZoneVadose NONE}, or no
\cmd{:SoilZone PHYSICAL}) cannot hold the columns' water: it warns, and \cmd{GWModelSummary.txt} reports the lines and their
water (above the columns' reference profiles) as dropped.
\end{itemize}
In Raven runs without groundwater, an age tracer (\cmd{:Transport AGE}; transport is not available in a coupled run,
Section~\ref{ch:inputs}) is restarted as follows. The file's tracer masses are volumes times entry times counted from a
datum: the start of the first run of a chain of restarts. \cmd{:StartTime} in \cmd{solution.rvc} is the start of the run that
wrote it; when the datum is earlier, the file also carries \cmd{:AgeTracerDatum} with that date. A run read from the file
counts its ages from \cmd{:AgeTracerDatum}, else from \cmd{:StartTime} when the file holds age-tracer masses; from a file
without them (a spin-up run without the tracer) the initial water has age 0.

Damaged groundwater values in the \cmd{.rvc} stop the run with a message naming the block: a value that is not a finite
number in any groundwater block, and negative water in \cmd{:GWSoilZoneColumns}, \cmd{:GWVadoseColumns} or
\cmd{:GWVadoseClassColumns}.

\subsection{Ensembles}
Raven's ensemble modes (Monte Carlo, DDS, EnKF) run the model several times in one program run. For every member the
groundwater model is rebuilt from scratch, MODFLOW is restarted, and its files go to that member's output folder.
Each member starts with a fresh water balance and, unless its own \cmd{.rvc} says
otherwise, from the routing memory and reservoir states of the first run (stage, outflow, loss terms). These defaults are
restored before the member's \cmd{.rvc} is read, so member-specific initial conditions (as used by EnKF forecasting) still
apply; groundwater hotstart values are likewise cleared, so each member uses only its own \cmd{.rvc}, and the steady-state
start is decided for each member separately
(Section~\ref{sec:coreflix}). An EnKF member's own time series file (\cmd{:ExtraRVTFilename}) may give
\cmd{:WellFlowRate} and \cmd{:BoundaryHead} series: they replace, for that member and the later ones, the series of the
same well or boundary read before, and a series naming no well or boundary stops the run.

\subsection{Safeguards}\label{sec:guard}
\begin{itemize}
\item \textbf{Non-converged steps.} Raven carries such a step on with MODFLOW's flows, as MODFLOW's own
\cmd{CONTINUE} option would: every flow MODFLOW booked is applied, and only its residual is not. A negative residual
(more outflow booked than the inflows and the storage release supply) is withheld from all discharge returned to Raven
--- seepage, river gains, discharge into wetland and lake stores, and with \cmd{:SoilZone PHYSICAL} the inflow to the
soil zone --- in proportion; what that discharge cannot cover stays in the aquifer. The withheld water is MODFLOW's
balance error of the step. Without \cmd{:SoilZone PHYSICAL}, the discharge columns of \cmd{GWBudget.csv} on such a step
(``seepage to Raven'', ``aquifer discharge to rivers'', ``aquifer discharge to surface stores'') give the discharge
handed to Raven, after the withholding, so each row's residual --- recharge plus storage release plus the other inflows,
less the discharge, less ``MF6 balance error'' --- is the volume withheld in that step; its run total is in
\cmd{GWModelSummary.txt} (``discharge withheld''). With \cmd{:SoilZone PHYSICAL} the withheld volume is booked per step
in its own columns of \cmd{GWBudget.csv} (seepage and soil inflow, river gains and surface-store gains withheld). The step is
flagged, and a warning written for the first ten such steps (test B8, Table~\ref{tab:vszinterface}).
\item \textbf{Rivers and stores} can never lose more water than they hold (Sections~\ref{ch:rivers},\ref{ch:stores}).
\item \textbf{Wells} cannot pump from a dry cell (\cmd{AUTO\_FLOW\_REDUCE}).
\item \textbf{Step count}: Raven stops if it would run past MODFLOW's last step.
\item \textbf{MODFLOW folder}: if Raven cannot enter the MODFLOW folder, the call fails with a clear message instead of
running MODFLOW in the wrong place.
\end{itemize}

